\documentclass[11pt,a4paper]{article}
\usepackage[margin=1.15in]{geometry}
\usepackage{amsmath,amssymb,amsthm}
\usepackage{booktabs}
\usepackage[hidelinks]{hyperref}

\theoremstyle{plain}
\newtheorem{theorem}{Theorem}
\newtheorem{proposition}[theorem]{Proposition}
\newtheorem{lemma}[theorem]{Lemma}
\newtheorem{corollary}[theorem]{Corollary}
\theoremstyle{definition}
\newtheorem{remark}[theorem]{Remark}

\newcommand{\R}{\mathbb{R}}
\newcommand{\Z}{\mathbb{Z}}
\newcommand{\T}{\mathbb{T}}
\newcommand{\Qodd}{\mathbb{Q}_{\mathrm{odd}}}
\newcommand{\al}{\alpha}
\newcommand{\ga}{\gamma}
\newcommand{\fr}[1]{\{#1\}}

\title{Accumulation points for the Collatz conjugacy map in $\mathbb{R}$}
\author{}
\date{}

\begin{document}
\maketitle

\begin{abstract}
The Collatz conjecture splits into two independent assertions: that no trajectory diverges, and
that the only cycle is the trivial one. Lagarias' periodicity conjecture is the divergence half,
extended from the integers to the rationals with odd denominator, and by the Bernstein--Lagarias
conjugacy it becomes a question about which aperiodic words can be parity vectors. L\'opez and
Stoll transported this question to $\R$, where a word with an accumulation point in its forced
trajectory is excluded as a witness; for their Example 6 they could not decide whether such a point
exists. We prove that it does. The argument rests on the identity $3^{\log2/\log3}=2$, which makes
every discontinuity of a certain limit function an upward jump and so produces an honest root of a
level set; the only computational input is a pair of inequalities verified in exact rational
interval arithmetic. We then embed the example in a one-parameter family, determine exactly for
which parameters one distinguished fibre is excluded, identify the limit of the family in closed
form, and isolate a scalar $R_n$ whose boundedness governs recurrence throughout the class. A
counterexample shows that recurrence is \emph{not} governed by the obvious candidates. We are
explicit about what this does not do: it does not bear on either half of the Collatz conjecture,
for reasons we set out precisely.
\end{abstract}

\section{Introduction}

\subsection*{The two halves}

Let $T$ be the $3x+1$ map, $T(x)=x/2$ for $x$ even and $T(x)=(3x+1)/2$ for $x$ odd. The Collatz
conjecture states that the trajectory of every positive integer reaches the cycle $1\mapsto2\mapsto1$.
It is the conjunction of two logically independent assertions:
\begin{itemize}
\item[(D)] \emph{no divergence}: every trajectory is bounded, hence eventually cyclic;
\item[(C)] \emph{no nontrivial cycle}: the only cycle among the positive integers is $\{1,2\}$.
\end{itemize}
Neither implies the other, and they are attacked by quite different means --- (C) by lower bounds on
the length of a hypothetical cycle, (D) by growth and density estimates.

The map $T$ extends to the $2$-adic integers $\Z_2$, and there it acts on the ring $\Qodd$ of
rationals with odd denominator, which is exactly the set of $x\in\Z_2$ whose orbit could be
rational. Lagarias' \emph{periodicity conjecture} asserts that the trajectory of every $x\in\Qodd$
is eventually cyclic. Restricted to the positive integers this is precisely (D); so the periodicity
conjecture \emph{is} the divergence half of Collatz, posed over a larger ring. It says nothing about
(C): eventual cyclicity does not assert which cycles occur. Thus
\[
   \text{periodicity conjecture}\ \Longrightarrow\ \text{(D)},
   \qquad
   \text{Collatz}=\text{(D)}\wedge\text{(C)} ,
\]
and a proof of the periodicity conjecture would settle the harder-looking half of Collatz while
leaving the cycle half untouched.

\subsection*{Words}

Bernstein and Lagarias \cite{BL} showed that the parity map $x\mapsto v(x)=(T^k(x)\bmod2)_{k\ge0}$
is a homeomorphism of $\Z_2$ conjugating $T$ to the shift, with inverse
\begin{equation}\label{eq:Phi}
   \Phi\Big(\sum_{i\ge0}2^{d_i}\Big)=-\sum_{i\ge0}\frac{2^{d_i}}{3^{i+1}},
   \qquad d_0<d_1<\cdots .
\end{equation}
Under this dictionary a trajectory is eventually cyclic exactly when its parity vector is eventually
periodic, so the periodicity conjecture becomes: \emph{no aperiodic word is the parity vector of an
element of $\Qodd$}. A counterexample would be a single aperiodic word $v$ with $\Phi(v)\in\Qodd$.
One is therefore led to examine aperiodic words one at a time and try to rule each out.

Two facts delimit the search. First \cite[Theorem 1]{LS}: if some $\zeta\in\Qodd$ has a divergent
trajectory, then necessarily $\liminf_\ell h(\ell)/\ell=\log2/\log3$, where $h(\ell)$ counts the
ones among the first $\ell$ digits of its parity vector. So a counterexample can only have this one
critical density, and every word considered below has it. Second \cite[p.~3]{LS}: an infinite subset
of $\Qodd$ has no accumulation point in $\R$; hence exhibiting one accumulation point in the orbit
of $\zeta$ proves $\zeta\notin\Qodd$ and eliminates that word.

\subsection*{The real picture}

L\'opez and Stoll \cite{LS} transported this to $\R$: for a word $v$ one reads \eqref{eq:Phi} as a
real series, obtaining $\Phi_\R(v)$ when it converges, and lets $T$ act on $\R$ through the branches
prescribed by $v$, irrespective of parity. Writing $\sigma$ for the shift, the \emph{forced
trajectory} of $\zeta=\Phi_\R(v)$ has terms $\Phi_\R(\sigma^tv)$. They eliminate every word of
supercritical density, and among critical-density words they settle several; for their Example 6
the question was left open:

\begin{quote}
``Example 6 is such a case where we cannot guarantee a (probably existing) accumulation point.''
\cite[p.~4]{LS}
\end{quote}

Our main result settles this case.

\begin{theorem}\label{thm:main}
Let $v$ be the word of \cite[\S14]{LS} and $\zeta=\Phi_\R(v)$. Then $\zeta$ is an accumulation point
of its own forced trajectory.
\end{theorem}

Sections \ref{sec:family}--\ref{sec:crit} develop the surrounding structure: the example lies in a
family indexed by a parameter $k$, for which the whole apparatus is parameter-free
(Proposition~\ref{prop:family}); one distinguished fibre of the reduction is excluded for exactly
$k\ge3$ (Theorem~\ref{thm:s0}); the family has an explicit limit (Proposition~\ref{prop:limit}); and
a single scalar governs recurrence throughout the class (Theorem~\ref{thm:Rn}), though not the one
the examples suggest (Remark~\ref{rem:counter}).

\subsection*{Where this sits}

It is worth being blunt about the distance to the conjectures. The programme sketched above would,
if carried out for \emph{every} aperiodic word, prove the periodicity conjecture and hence half (D)
of Collatz; half (C) would remain untouched by any of it. What is proved here is one step of that
programme --- a single word, and then a family --- together with two structural facts about how far
the step generalises.

Both facts are negative, and both are worth recording. The associated word of
\cite[Example 7]{LS} has the same critical density and provably \emph{no} accumulation point, so the
method cannot be universal: it will never eliminate all aperiodic words. And
Corollary~\ref{cor:irr} constrains the \emph{real} value $\Phi_\R(v)$, whereas the periodicity
conjecture concerns the \emph{$2$-adic} value $\Phi(v)$; these are limits of the same rational
partial sums in different completions, and we cannot show they agree. Section~\ref{sec:scope}
locates that missing implication exactly, as a growth estimate \eqref{eq:bridge} for the orbit of
the hypothetical rational, which is the sharpest form in which we can state what is absent. Nothing
below should be read as progress on the Collatz conjecture itself.

\section{Notation}\label{sec:setup}

Throughout $\al=\log2/\log3$ and $\ga=\al/(1-\al)=\log2/\log(3/2)$, so that
\begin{equation}\label{eq:key}
  3^{\al}=2,\qquad \ga\log\tfrac32=\log2 .
\end{equation}
Every constant below is a consequence of \eqref{eq:key}. We write $\fr{x}$ for the fractional part.

Fix $k\ge2$ and set $f(n)=\lceil n/k\rceil$, $g(n)=n-f(n)$. The word is
\begin{equation}\label{eq:word}
  v=\prod_{n\ge1}1^{m_n}0^{f(n)},\qquad
  m_n=\left\lfloor\frac{(L_{n-1}+n)\al-H_{n-1}}{1-\al}\right\rfloor ,
\end{equation}
where $L_n,H_n$ are the length and the number of ones of the first $n$ factors, $L_0=H_0=0$; thus
$L_n=L_{n-1}+m_n+f(n)$ and $H_n=H_{n-1}+m_n$. Example 6 of \cite{LS} is $k=6$; its first values are
$m_1,\dots,m_8=1,4,3,3,4,3,4,5$. Let $Z_n$ be the number of zeros among the first $n$ factors and
$G(n)=Z_{n-1}+n$; from $f(n)=\lceil n/k\rceil$ one computes that $G(kM+s)$ is a quadratic polynomial
in $M$ with leading coefficient $k/2$.

Put $D_n=\al L_n-H_n$. Let $z_n=\Phi_\R(\sigma^{L_n}v)$ be the value of the forced trajectory at the
end of the $n$-th factor, so $z_0=\zeta$.

\begin{lemma}\label{lem:drift}
$H_n=\lfloor\ga G(n)\rfloor$ and $D_n=-\al g(n)+\beta_n$ with $\beta_n=(1-\al)\fr{\ga G(n)}\in[0,1-\al)$.
Consequently
\begin{equation}\label{eq:ramp}
   D_{n+r}-D_n\le-\al\Big\lfloor\tfrac{(k-1)r}{k}\Big\rfloor+(1-\al),
   \qquad
   D_n\le(1-\al)-\tfrac{(k-1)\al}{k}(n-1).
\end{equation}
\end{lemma}

\begin{proof}
As $H_{n-1}\in\Z$ it may be absorbed into the floor in \eqref{eq:word}, giving
\[
  H_n=\left\lfloor\frac{\al L_{n-1}+n\al-\al H_{n-1}}{1-\al}\right\rfloor
     =\big\lfloor\ga(L_{n-1}-H_{n-1}+n)\big\rfloor=\lfloor\ga G(n)\rfloor .
\]
Then $D_n=\al Z_n-(1-\al)H_n=\al(Z_n-G(n))+\beta_n$ and $Z_n-G(n)=f(n)-n=-g(n)$. The estimates
follow from $\beta_{n+r}-\beta_n<1-\al$ and $g(n)=\lfloor(k-1)n/k\rfloor$.
\end{proof}

\section{Proof of Theorem \ref{thm:main}}\label{sec:proof}

Throughout this section $k=6$. Let $d_0<d_1<\cdots$ be the positions of the ones of $v$.

\begin{lemma}\label{lem:stage}
$\displaystyle z_n=-\sum_{r\ge0}3^{\,D_{n+r}-D_n}\Big(1-(2/3)^{m_{n+r+1}}\Big)$.
\end{lemma}

\begin{proof}
The ones of $\sigma^{L_n}v$ are the $d_{H_n+i}-L_n$, so
$z_n=-\sum_{i\ge0}2^{\,d_{H_n+i}-L_n}3^{-(i+1)}$. The ones of the factor $n+r+1$ occupy the
positions $L_{n+r},\dots,L_{n+r}+m_{n+r+1}-1$ and carry the indices
$H_{n+r},\dots,H_{n+r}+m_{n+r+1}-1$, so they contribute
$-2^{\,L_{n+r}-L_n}3^{-(H_{n+r}-H_n)}\cdot\tfrac13\sum_{i<m}(2/3)^i$ with $m=m_{n+r+1}$. The inner
sum is $3(1-(2/3)^m)$, and $2^{\Delta L}3^{-\Delta H}=3^{\,\al\Delta L-\Delta H}=3^{\,D_{n+r}-D_n}$
by \eqref{eq:key}.
\end{proof}

Let $w_n=-\sum_{r\ge0}3^{\,D_{n+r}-D_n}$ be the corresponding sum with all damping factors removed.

\begin{proposition}\label{prop:sat}
$0\le z_n-w_n\le 9\rho^{\,n+1}$ with $\rho=2^{-1/6}$.
\end{proposition}

\begin{proof}
By \eqref{eq:ramp} and \eqref{eq:key}, $3^{\,D_{n+r}-D_n}\le3^{1-\al}3^{\al}q_1^{\,r}=3q_1^{\,r}$
where
\begin{equation}\label{eq:q1}
  q_1=3^{-5\al/6}=2^{-5/6}.
\end{equation}
From \eqref{eq:word}, $m_n\ge\ga(n+5)/6-1$, and $(2/3)^{\ga/6}=2^{-1/6}=\rho$ by \eqref{eq:key}, so
$(2/3)^{m_{n+r+1}}\le\tfrac32\rho^{\,n+r+1}$. The $r$-th term of
$z_n-w_n=\sum_r3^{\,D_{n+r}-D_n}(2/3)^{m_{n+r+1}}$ is therefore at most
$\tfrac92q_1^{\,r}\rho^{\,n+r+1}$. Now
\[
   q_1\rho=2^{-5/6}\cdot2^{-1/6}=\tfrac12 ,
\]
so the series is geometric of ratio $\tfrac12$ and sums to at most $9\rho^{\,n+1}$.
\end{proof}

The coincidence $q_1\rho=\tfrac12$ is not accidental: by \eqref{eq:key} the loss of drift and the
gain of damping per factor are reciprocal at the critical density, for every $k$.

Write $n=6M+s$ with $0\le s\le5$, and set $u=\fr{\ga G(n)}$, $v=\fr{\ga M}$ (no confusion with the
word should arise). Since $6M\equiv0\pmod 6$ we have
$f(n+r)-f(n)=\lceil(s+r)/6\rceil-\lceil s/6\rceil$, and expanding the quadratic $G$ gives
$G(n+r)-G(n)=rM+C_s(r)$ with $C_s(r)=G(s+r)-G(s)$ free of $M$. As $r\lfloor\ga M\rfloor\in\Z$,
$\fr{\ga G(n+r)}=\fr{u+rv+\ga C_s(r)}$. With
$A_s(r)=\al\big((\lceil(s+r)/6\rceil-\lceil s/6\rceil)-r\big)$, Lemma~\ref{lem:drift} therefore
yields
\begin{equation}\label{eq:F}
  w_{6M+s}=F_s(u,v),\qquad
  F_s(u,v)=-\sum_{r\ge0}3^{\,A_s(r)+(1-\al)(\fr{u+rv+\ga C_s(r)}-u)} .
\end{equation}
Thus $w$ is a fixed function of two variables sampled along an orbit. Since
$A_s(r)\le\al(1-5r/6)$ and the second exponent is less than $1-\al$, every term of \eqref{eq:F} is
smaller than $3q_1^{\,r}$; the series converges uniformly and $F_s$ has one-sided limits everywhere.
Moreover $\fr\cdot$ is right-continuous, hence so is $F_s$.

\begin{lemma}\label{lem:up}
Every discontinuity of $F_s$ is an upward jump.
\end{lemma}

\begin{proof}
As a phase $u+rv+\ga C_s(r)$ crosses an integer upward, $\fr\cdot$ drops by $1$, so the exponent of
that term drops by $1-\al$ and the term is multiplied by $3^{\al-1}=2/3$ by \eqref{eq:key}. Every
term is positive, so each jump of $F_s=-\sum$ is upward; uniform convergence transfers this to the
sum.
\end{proof}

\begin{lemma}\label{lem:cross}
Let $f:[a,b]\to\R$ have one-sided limits everywhere, be right-continuous, and satisfy
$f(x^-)\le f(x)$ for all $x$. If $f(a)>c>f(b)$ then $x^*=\sup\{x:f(x)>c\}$ satisfies $f(x^*)=c$, and
$f$ has no jump at $x^*$.
\end{lemma}

\begin{proof}
The set $A=\{x:f(x)>c\}$ contains $a$. If $x^*=b$ then $x_j\uparrow b$ with $f(x_j)>c$ would give
$f(b)\ge f(b^-)\ge c$, contradicting $f(b)<c$; so $x^*<b$. No $x>x^*$ lies in $A$, so $f\le c$ on
$(x^*,b]$ and right-continuity gives $f(x^*)\le c$. Points of $A$ increasing to $x^*$ give
$f(x^{*-})\ge c$, and $f(x^*)\ge f(x^{*-})$. Hence $c\le f(x^{*-})\le f(x^*)\le c$.
\end{proof}

\begin{lemma}\label{lem:2d}
If $F_s(u_0,\cdot)$ has no jump at $v^*$ and $u_0\in(0,1)$, then $F_s$ is continuous at $(u_0,v^*)$
as a function on $\T^2$.
\end{lemma}

\begin{proof}
The jump of $F_s(u_0,\cdot)$ at $v^*$ is the sum of the jumps of the individual terms, all of one
sign by Lemma~\ref{lem:up}; a sum of nonnegative terms vanishes only if each does. For $r\ge1$ the
phase is strictly monotone in $v$, so no phase line $u+rv+\ga C_s(r)\in\Z$ passes through
$(u_0,v^*)$; the term $r=0$ is constant. Now given $\varepsilon>0$ choose $R$ with
$3q_1^{\,R+1}/(1-q_1)<\varepsilon/3$. The tail contributes less than $\varepsilon/3$ at every point
of $\T^2$, while the head is a finite sum whose discontinuity set is a finite union of lines, none
through $(u_0,v^*)$; it therefore varies by less than $\varepsilon/3$ on a ball there.
\end{proof}

The care in Lemma~\ref{lem:2d} is necessary: the full discontinuity set is a countable union of
lines of spacing $1/r$, hence dense, and no jump-free neighbourhood exists.

\begin{lemma}\label{lem:weyl}
For each $s$, the sequence $\big(\fr{\ga G(6M+s)},\fr{\ga M}\big)_{M\ge0}$ is equidistributed in
$\T^2$.
\end{lemma}

\begin{proof}
$\ga$ is irrational, since $\ga=p/q$ would give $2^{q+p}=3^p$. For $(h_1,h_2)\ne(0,0)$ the
polynomial $h_1\ga G(6M+s)+h_2\ga M$ has an irrational non-constant coefficient, namely $3h_1\ga$ if
$h_1\ne0$ and $h_2\ga$ otherwise. Apply Weyl's theorem and the Weyl criterion on $\T^2$.
\end{proof}

\begin{proposition}\label{prop:straddle}
$F_3\big(\tfrac{37}{100},\tfrac{369}{4000}\big)>\zeta>F_3\big(\tfrac{37}{100},\tfrac{37}{400}\big)$.
\end{proposition}

This is a finite computation, carried out in interval arithmetic whose endpoints are exact
rationals: $\log2$ and $\log3$ come from the $\operatorname{artanh}$ series with its explicit
remainder, $3^x$ from an argument-reduced exponential with the Lagrange remainder, $\zeta$ from its
defining series with the geometric tail of \eqref{eq:ramp}, and $F_3$ from \eqref{eq:F} truncated at
$r=150$ with the tail bound $3q_1^{151}/(1-q_1)<9.1\cdot10^{-38}$. Each floor and fractional part
occurring is certified not to straddle an integer. The certified values are
\[
\begin{array}{ll}
  \zeta &=-2.2923055020547008507324073557\ldots\\
  F_3(\tfrac{37}{100},\tfrac{369}{4000}) &=-2.2922631366352524359395122553\ldots\\
  F_3(\tfrac{37}{100},\tfrac{37}{400}) &=-2.2925637964328159779770742733\ldots
\end{array}
\]
with margins $4.2\cdot10^{-5}$ and $2.6\cdot10^{-4}$ against arithmetic widths below $10^{-37}$.

\begin{proof}[Proof of Theorem \ref{thm:main}]
Take $s=3$ and $u_0=37/100$. By Proposition~\ref{prop:straddle} and Lemmas~\ref{lem:up},
\ref{lem:cross}, the function $F_3(u_0,\cdot)$ meets the level $\zeta$ at
$v^*=\sup\{v:F_3(u_0,v)>\zeta\}$ without a jump there, so $F_3(u_0,v^*)=\zeta$ and, by
Lemma~\ref{lem:2d}, $F_3$ is continuous at $(u_0,v^*)$. By Lemma~\ref{lem:weyl} there are
$M_1<M_2<\cdots$ with $(\fr{\ga G(6M_j+3)},\fr{\ga M_j})\to(u_0,v^*)$, whence
$w_{6M_j+3}=F_3(\cdot)\to\zeta$ by \eqref{eq:F}, and $z_{6M_j+3}\to\zeta$ by
Proposition~\ref{prop:sat}.
\end{proof}

\begin{remark}
The conclusion is a genuine one about the trajectory, not only about the auxiliary sequence:
computing $z_{6M+3}$ directly from Lemma~\ref{lem:stage} at orbit points approaching $(u_0,v^*)$
gives $|z_{6M+3}-\zeta|$ proportional to the distance, with ratio between $0.88$ and $1.20$ at the
five closest points with $M<2\cdot10^7$ (the nearest being $M=4\,153\,063$ at distance
$9.4\cdot10^{-5}$).
\end{remark}

\section{A family}\label{sec:family}

Nothing in \S\ref{sec:proof} used the value $6$ except through $f$. Keep \eqref{eq:word} with
$f(n)=\lceil n/k\rceil$ and write $\zeta_k$, $F^{(k)}_s$ for the resulting objects.

\begin{proposition}\label{prop:family}
For every $k\ge2$, Lemmas~\ref{lem:drift}--\ref{lem:weyl} and Proposition~\ref{prop:sat} hold with
\[
   q_1=3^{-\al(k-1)/k}=2^{-(k-1)/k},\qquad \rho=(2/3)^{\ga/k}=2^{-1/k},
   \qquad q_1\rho=\tfrac12 .
\]
The ratios $3^{\al-1}=2/3$, $3^{1-\al}=3/2$ and $3^{-\al}=1/2$ do not involve $k$, so
Lemma~\ref{lem:up} is unchanged; $C_s(r)$ is free of $M$; and $G(kM+s)$ is quadratic in $M$ with
leading coefficient $k/2\ne0$.
\end{proposition}

Consequently the proof of Theorem~\ref{thm:main} applies verbatim once a level crossing is
exhibited.

\begin{theorem}\label{thm:family}
Let $k\ge2$. If $F^{(k)}_s(u_0,v_1)>\zeta_k>F^{(k)}_s(u_0,v_2)$ for some $s$ and $v_1<v_2$, then
$\zeta_k$ is an accumulation point of its own forced trajectory.
\end{theorem}

The hypothesis has been certified, by the computation described after
Proposition~\ref{prop:straddle}, in the five cases of Table~\ref{tab:family}.

\begin{table}[h]
\centering
\begin{tabular}{rrlll}
\toprule
$k$ & $s$ & $u_0$ & $[v_1,v_2]$ & $\zeta_k$\\
\midrule
$6$ & $3$ & $37/100$ & $[\tfrac{369}{4000},\tfrac{37}{400}]$ & $-2.29230550205470085\ldots$\\
$4$ & $2$ & $1/60$   & $[\tfrac{257}{400},\tfrac{129}{200}]$ & $-2.43242829803240526\ldots$\\
$5$ & $1$ & $7/40$   & $[\tfrac{9}{25},\tfrac{73}{200}]$     & $-2.34442124302516940\ldots$\\
$8$ & $3$ & $1/40$   & $[\tfrac{4}{25},\tfrac{33}{200}]$     & $-2.26316625669379821\ldots$\\
$2$ & $0$ & $19/24$  & $[\tfrac{47}{75},\tfrac{377}{600}]$   & $-3.51194122554282436\ldots$\\
\bottomrule
\end{tabular}
\caption{Certified crossings. Margins exceed $4\cdot10^{-5}$ throughout, against arithmetic widths
below $10^{-21}$.}\label{tab:family}
\end{table}

\section{The fibre \texorpdfstring{$s=0$}{s=0}}\label{sec:s0}

For $k\ge3$ the crossings above occur only for $s\ge1$, and this is not an accident of the search.
Write $\lambda=3^{1-\al}=3/2$ and $a_r=3^{A_0(r)}=2^{\lceil r/k\rceil-r}$, so that
$-F^{(k)}_0(u,v)=1+\sum_{r\ge1}a_r\lambda^{\,x_r-u}$ with $x_r$ the relevant fractional part.

\begin{proposition}\label{prop:s0}
$-F^{(k)}_0(u,v)>B_k$ for all $(u,v)$, where
\[
   B_k=1+\tfrac23\sum_{r\ge1}2^{\lceil r/k\rceil-r}=1+\frac43\cdot\frac{1-2^{-k}}{1-2^{-(k-1)}} .
\]
$B_k$ decreases strictly to $7/3$, so $B_k>7/3$ for every $k$.
\end{proposition}

\begin{proof}
The term $r=0$ equals $1$; for $r\ge1$ the factor $\lambda^{\,x_r-u}$ exceeds
$\lambda^{-1}=\tfrac23$, and $3^{A_0(r)}=2^{\lceil r/k\rceil-r}$ by \eqref{eq:key}. Writing
$r=(b-1)k+i$ with $1\le i\le k$ gives
$\sum_{r\ge1}2^{\lceil r/k\rceil-r}=(1-2^{-k})\cdot2/(1-2^{-(k-1)})$. Finally
$x\mapsto(1-x)/(1-2x)$ is increasing and $x=2^{-k}$ decreases.
\end{proof}

$B_k$ suffices for $k\ge5$ but not below, and the deficiency is repaired by observing that the first
two terms are coupled: $C_0(1)=1$ and $C_0(2)=3$ for every $k$, so with $x=\fr{u+v+\ga}$ one has
$x_2=\fr{2x+\ga-u}$.

\begin{proposition}\label{prop:two}
$\displaystyle\inf_{u,v}\Big(\lambda^{\,x-u}+\tfrac12\lambda^{\,x_2-u}\Big)=\frac1{\sqrt2}+\frac13$,
whence $-F^{(k)}_0>B_k^{\sharp}:=B_k+\tfrac1{\sqrt2}-\tfrac23$.
\end{proposition}

\begin{proof}
Put $e=x-u$ and $m=\lfloor2e+u+\ga\rfloor\in\{0,1,2,3\}$, so that $x_2-u=2e+\ga-m$ and the quantity
is $\lambda^{\,e}+\tfrac12\lambda^{\,2e+\ga-m}$, strictly increasing in $e$. On the branch $m$ the
constraint $m\le2e+u+\ga$ with $u<1$ gives $e>(m-1-\ga)/2$, and $x\ge0$ gives $e>-1$. Evaluating at
these endpoints gives $1.1112\ldots$, $\tfrac1{\sqrt2}+\tfrac13$, $1.1993\ldots$, $1.3936\ldots$ for
$m=0,1,2,3$. The minimum is at $m=1$, $e=-\ga/2$, where $\lambda^{-\ga/2}=(3/2)^{-\ga/2}=2^{-1/2}$
by \eqref{eq:key} and $\tfrac12\lambda^{-1}=\tfrac13$. The two terms were assigned
$\tfrac23(1+\tfrac12)=1$ in $B_k$.
\end{proof}

For $k=3$ even $B_3^{\sharp}$ falls short, and one uses all of the coupling.

\begin{proposition}\label{prop:QR}
Let $\theta_r=rv+\ga C_0(r)$ and $Q_R=\inf_v\sum_{r\le R}a_r\lambda^{\fr{\theta_r}}$. Then
$-F^{(k)}_0\ge B_k+\tfrac23\big(Q_R-\sum_{r\le R}a_r\big)$, and $Q_R$ is the minimum of the finitely
many values at $v=(n-\ga C_0(r'))/r'$, $1\le r'\le R$, $0\le n<r'$.
\end{proposition}

\begin{proof}
Since $x_r-u=\fr{\theta_r}-[\,u\ge1-\fr{\theta_r}\,]$ and $\lambda>1$, each term is at least
$\tfrac23a_r\lambda^{\fr{\theta_r}}$ for every $u$; bounding the terms $r>R$ by $\tfrac23a_r$ gives
the inequality. On a cell where all $\lfloor\theta_r\rfloor$ are constant, $\fr{\theta_r}$ is linear
in $v$ with slope $r>0$, so every summand increases and the infimum over the cell is at its left
endpoint. These are the stated points, taken modulo $1$, the configuration being $1$-periodic. At
such a point $\theta_r=p+\ga c$ with $p,c$ rational, and as $\ga$ is irrational this is an integer
only when $c=0$ and $p\in\Z$, a case decided exactly.
\end{proof}

\begin{theorem}\label{thm:s0}
For $k\ge3$ one has $F^{(k)}_0<\zeta_k$ everywhere, so the hypothesis of
Theorem~\ref{thm:family} cannot hold with $s=0$. For $k=2$ it does hold with $s=0$
(Table~\ref{tab:family}).
\end{theorem}

\begin{proof}
For $k\ge7$: the words for different $k$ agree on their first $k$ factors, since
$\lceil n/k\rceil=1$ for $n\le k$; so the sum of the first seven terms of
$|\zeta_k|=\sum_n3^{D_{n-1}}(1-(2/3)^{m_n})$ is independent of $k$, and \eqref{eq:ramp} bounds the
remainder by $\tfrac32q_1^{\,6}/(1-q_1)$ with $q_1\le2^{-6/7}$. This gives
$|\zeta_k|\le2.3231<\tfrac73<B_k$. For $k=6,5$ one has $|\zeta_6|=2.29231<B_6=2.35484$ and
$|\zeta_5|=2.34442<B_5=2.37778$; note $|\zeta_5|>7/3$, which is why $B_k$ rather than its limit is
used. For $k=4$, $B_4=17/7$ is too small but
$B_4^{\sharp}=\tfrac{17}7+\tfrac1{\sqrt2}-\tfrac23=2.46901>2.43243=|\zeta_4|$. For $k=3$,
Proposition~\ref{prop:QR} with $R=8$ gives $-F^{(3)}_0\ge2.69267>2.66441=|\zeta_3|$.
\end{proof}

At $k=2$ every version of the bound fails, as it must; running the $R=8$ estimate there returns
$3.189<3.512=|\zeta_2|$ and correctly declines to conclude.

\begin{proposition}\label{prop:limit}
As $k\to\infty$ the words agree on ever longer initial segments and converge to the word with one
zero per factor and $H_n=\lfloor\ga(2n-1)\rfloor$, $H_0=0$, $L_n=H_n+n$. Correspondingly
$\zeta_k\to\zeta_\infty$ with
\[
   |\zeta_\infty|=1+\frac12\sum_{n\ge1}2^{\,n}\Big(\frac23\Big)^{\lfloor\ga(2n-1)\rfloor}
   =2.2548974992816421641\ldots
\]
\end{proposition}

\begin{proof}
With $f\equiv1$ one has $Z_n=n$ and $G(n)=2n-1$, so Lemma~\ref{lem:drift} gives $H_n$ and $L_n$. The
$n$-th factor contributes $2^{L_{n-1}}3^{-H_{n-1}}(1-(2/3)^{m_n})=2^{\,n-1}A_{n-1}(1-A_n/A_{n-1})$
to $|\zeta_\infty|$, where $A_n=(2/3)^{H_n}$; summing $\sum_n2^{\,n-1}(A_{n-1}-A_n)$ by parts gives
$A_0+\tfrac12\sum_n2^{\,n}A_n$. Convergence is geometric of ratio $\tfrac12$ by \eqref{eq:key}.
\end{proof}

The limit word has ones of density $2\ga/(2\ga+1)$, not $\al$, so it is not itself a member of the
family; Proposition~\ref{prop:limit} concerns the limit of the values.

\section{A criterion for recurrence}\label{sec:crit}

The examples above raise the question of which words of the shape \eqref{eq:word} have recurrent
trajectories at all. By Lemma~\ref{lem:drift},
\begin{equation}\label{eq:split}
   3^{\,D_{n+r}-D_n}=2^{-(g(n+r)-g(n))}\cdot3^{\,\beta_{n+r}-\beta_n},
   \qquad \tfrac23<3^{\,\beta_{n+r}-\beta_n}<\tfrac32 ,
\end{equation}
splitting the ratio into a deterministic part and a bounded fluctuation. This suggests the scalar
\[
   R_n=\sum_{r\ge0}2^{-(g(n+r)-g(n))}.
\]

\begin{theorem}\label{thm:Rn}
Let $f$ in \eqref{eq:word} be nondecreasing with $f(n)\to\infty$, the density being critical, and
suppose $\Phi_\R(v)$ exists. Then $\tfrac29R_n\le-z_n\le\tfrac32R_n$ for all large $n$, so
$(z_n)$ has an accumulation point if and only if $\liminf_nR_n<\infty$. Equivalently, if a run of
$b_j$ zeros begins at time $t_j$ then $R_{t_j+b_j}=2^{-b_j}R_{t_j}$, and recurrence holds iff
$\liminf_j2^{-b_j}R_{t_j}<\infty$.
\end{theorem}

The constants come from \eqref{eq:split} and from the damping factor of Lemma~\ref{lem:stage} being
eventually at least $\tfrac13$. A finite $\liminf$ gives a bounded infinite subsequence of the
$z_n$ and hence a convergent one; if $R_n\to\infty$ then the local minima of $R$, which occur at the
ends of the zero runs, escape, and with them the trajectory.

\begin{corollary}\label{cor:irr}
Under the hypotheses of Theorem~\ref{thm:Rn}, $\liminf_nR_n<\infty$ implies $\Phi_\R(v)\notin\Qodd$.
\end{corollary}

Indeed an infinite subset of $\Qodd$ has no accumulation point in $\R$ \cite[p.~3]{LS}. For the word
of Theorem~\ref{thm:main}, $g(n)=\lfloor5n/6\rfloor$ gives $R_n\to(1-2^{-5/6})^{-1}=2.2791\ldots$,
so \S\ref{sec:proof} is the bounded case of Theorem~\ref{thm:Rn}.

\begin{remark}\label{rem:counter}
Neither obvious candidate governs recurrence. Zero runs tending to infinity do not suffice: taking
$g(n)=\lceil c\log(n+1)\rceil$ in \eqref{eq:word}, so that $f(n)=n-g(n)\sim n$ and the zero runs grow
\emph{linearly}, one still has critical density, yet $R_n$ diverges and no accumulation point exists.
What matters is whether $g$ outruns $\log_2$, since $R_n$ is bounded exactly when
$g(n+r)-g(n)\gtrsim(1+\varepsilon)\log_2r$. Nor is the drift the invariant: $D_n\to-\infty$ governs
convergence of the series, not recurrence, and the word of Theorem~\ref{thm:main} and its associated
word \cite[Example 7]{LS} both have $D_n\to-\infty$ with opposite behaviour. By the zero-run form of
Theorem~\ref{thm:Rn} the relevant quantity is not the contraction length $b_j$ alone but the scale
$R_{t_j}$ on which it acts.
\end{remark}

\section{Scope}\label{sec:scope}

Theorem~\ref{thm:main} settles the case left open in \cite{LS} and, by \cite[p.~3]{LS}, gives
$\zeta\notin\Qodd$ for that word. It bears on neither half of the Collatz conjecture. Half (C), the
absence of nontrivial cycles, is untouched by everything here: cycles correspond to eventually
periodic words, and every word considered is aperiodic by construction. Half (D) would follow from
the periodicity conjecture, but our results fall short of that for two reasons worth stating
precisely.

First, the conjecture quantifies over all aperiodic words, and no uniform mechanism is available
here: the associated word of \cite[Example 7]{LS} has the same critical density and provably no
accumulation point, while the method above is parasitic on the structure of \eqref{eq:word} --- $G$
quadratic and $C_s(r)$ free of $M$ --- which a general word does not possess. Theorem~\ref{thm:Rn}
identifies when recurrence holds in the class \eqref{eq:word}, but it yields only the existence of
\emph{some} accumulation point, not that $\zeta$ is one of them; the latter needs the recurrence of
the reduction to a phase at which \eqref{eq:F} takes the value $\zeta$, which is what
Lemma~\ref{lem:weyl} supplies here and what a general argument would have to replace.

Second, Corollary~\ref{cor:irr} constrains the real value $\Phi_\R(v)$, whereas the conjecture
concerns the $2$-adic value $\Phi(v)$. These are limits of the same rational partial sums
$S_\ell=\Phi_\R(u^{(\ell)})$ in different completions and need not agree. The missing implication is
that $\Phi(v)\in\Qodd$ forces $\Phi_\R(v)=\Phi(v)$ when the real series converges. It can be located
exactly: for a finite prefix there is no ambiguity, $\Phi(u)=\Phi_\R(u)$, and combining
$T^\ell(x)=(\varphi(u)+3^hx)/2^\ell$ with $\Phi(u)=-\varphi(u)/3^h$ gives
\begin{equation}\label{eq:bridge}
   x-S_\ell=\frac{2^\ell}{3^h}\,T^\ell(x)=3^{\,D_\ell}\,T^\ell(x),
\end{equation}
so that $\Phi_\R(v)=x$ if and only if $T^\ell(x)=o(3^{-D_\ell})$, while $S_\ell\to L\ne x$ forces
$T^\ell(x)\sim3^{-D_\ell}(x-L)$. This is a relocation rather than a weakening: $3^{-D_\ell}$ is the
natural homogeneous growth scale of the maps being iterated, so the condition carries no slack and
is precisely the cancellation selecting $x=\Phi_\R(v)$. What \eqref{eq:bridge} provides is a
concrete form for what is missing --- a growth estimate for the orbit of the hypothetical rational
--- rather than an abstract incompatibility of completions. We have not been able to supply it.

Finally, Proposition~\ref{prop:straddle} and Table~\ref{tab:family} are computer-assisted, and their
correctness depends on that of the interval arithmetic; the code is short, uses no third-party
library, and was checked against an independent implementation agreeing to twenty significant
digits. We have not found the scalar $R_n$ or the equivalence of Theorem~\ref{thm:Rn} in \cite{LS},
but make no claim of novelty beyond that comparison.

\begin{thebibliography}{9}
\bibitem{BL} D.~J.~Bernstein and J.~C.~Lagarias, \emph{The $3x+1$ conjugacy map},
Canad. J. Math. \textbf{48} (1996), 1154--1169.
\bibitem{LS} J.~L\'opez and P.~Stoll, \emph{The $3x+1$ periodicity conjecture in $\mathbb{R}$},
arXiv:2101.12747 (2021).
\bibitem{MY} K.~G.~Monks and J.~Yazinski, \emph{The autoconjugacy of the $3x+1$ function},
Discrete Math. \textbf{275} (2004), 219--236.
\bibitem{KN} L.~Kuipers and H.~Niederreiter, \emph{Uniform Distribution of Sequences}, Wiley, 1974.
\end{thebibliography}

\end{document}
